% Generated cross-document labels pointing to main.pdf.
% Refresh with build_pdfs.sh after structural edits.
\newlabel{sec:intro}{{I}{1}{Introduction}{section.1}{main.pdf}}
\newlabel{sec:motivation}{{II}{2}{A Motivating Example and Our Approach}{section.2}{main.pdf}}
\newlabel{fig:example}{{1}{3}{Motivating Example}{figure.1}{main.pdf}}
\newlabel{fig:overallarchitecture}{{2}{3}{Overview of Our Approach}{figure.2}{main.pdf}}
\newlabel{sec:summarygeneration}{{III}{4}{Pre-/Postcondition Generation}{section.3}{main.pdf}}
\newlabel{sec:default-pre}{{\mbox  {III-B}}{5}{Default Precondition Generation}{subsection.3.2}{main.pdf}}
\newlabel{sec:functionsummary}{{\mbox  {III-C}}{5}{Function Specification Generation Algorithm}{subsection.3.3}{main.pdf}}
\newlabel{alg:funspec-gen}{{1}{5}{Function Specification Generation Algorithm}{algocf.1}{main.pdf}}
\newlabel{sec:llm-fallback}{{\mbox  {III-D}}{6}{LLM Fallback}{subsection.3.4}{main.pdf}}
\newlabel{sec:loopinvariant}{{IV}{6}{Loop Invariant Generation}{section.4}{main.pdf}}
\newlabel{Def:invariant}{{2}{6}{}{definition.2}{main.pdf}}
\newlabel{sec:template}{{\mbox  {IV-B}}{6}{Template Generation}{subsection.4.2}{main.pdf}}
\newlabel{sec:refinement}{{\mbox  {IV-C}}{7}{Verifier-Guided Specification Refinement}{subsection.4.3}{main.pdf}}
\newlabel{sec:experiments}{{V}{7}{Experiments}{section.5}{main.pdf}}
\newlabel{tab:dataset_overview}{{I}{8}{Dataset overview by program type}{table.1}{main.pdf}}
\newlabel{tab:performance_rates_model}{{II}{8}{Specification generation performance under Pass@1 and Pass@5. The complete Pass@1/3/5 table is given in Table~\supref {tab:performance_rates_model_full} (\appref {app:rq2-full})}{table.2}{main.pdf}}
\newlabel{tab:ablation}{{III}{9}{Ablation study}{table.3}{main.pdf}}
\newlabel{sec:inv_experiments}{{\mbox  {V-E}}{9}{RQ4: Does \toolname \ outperform LLM- and SMT-based baselines?}{subsection.5.5}{main.pdf}}
\newlabel{sec:large-scale-experiment}{{VI}{9}{Goal-Independent Evaluation}{section.6}{main.pdf}}
\newlabel{subsec:rq5-valid}{{\mbox  {VI-B}}{9}{RQ5: How often does \toolname \ produce a \emph {Valid} specification?}{subsection.6.2}{main.pdf}}
\newlabel{tab:comparison}{{IV}{10}{Accuracy comparison across benchmarks and baselines. \toolname \ reports Pass@5 (cf.\ Table~\ref {tab:performance_rates_model}); baselines use their best reported configuration}{table.4}{main.pdf}}
\newlabel{tab:large_experiment_summary}{{V}{10}{Native-verifier validity on \toolname -400}{table.5}{main.pdf}}
\newlabel{subsec:rq6-strength}{{\mbox  {VI-C}}{10}{RQ6: How strong are \toolname 's specifications versus the baselines?}{subsection.6.3}{main.pdf}}
\newlabel{fig:fc_vs_judge_strength}{{3}{11}{Dimension-wise relations for \autospec \ versus \toolname \ over identical data sets. A directional segment means a more permissive accepted domain for preconditions and a logically stronger predicate for postconditions or invariants. ``Judge + expert'' is the corrected judge result; WP checks the same implications. Right-hand counts follow the legend order}{figure.3}{main.pdf}}
\newlabel{subsec:rq7-cost}{{\mbox  {VI-D}}{11}{RQ7: What does \toolname \ cost in time, tokens, API calls, and memory?}{subsection.6.4}{main.pdf}}
\newlabel{subsec:rq8-failure}{{\mbox  {VI-E}}{11}{RQ8: What are \toolname 's dominant failure modes?}{subsection.6.5}{main.pdf}}
\newlabel{sec:relatedwork}{{VII}{11}{Related Work}{section.7}{main.pdf}}
\newlabel{fig:rq7-cost}{{4}{12}{Mean cost per attempted program on all 400 \texttt {gpt-5.4-mini} cases. Each bar stacks valid-run (solid) and failed-run (hatched) contributions, each divided by 400}{figure.4}{main.pdf}}
\newlabel{fig:failure_root_causes}{{5}{12}{Failure outcome and root-cause summary for the \texttt {gpt-5.4-mini} \toolname -400 rerun}{figure.5}{main.pdf}}
\newlabel{sec:conclusion}{{VIII}{12}{Conclusion}{section.8}{main.pdf}}
\newlabel{tab:positioning}{{VI}{13}{Positioning of \toolname \ against representative specification and invariant generation tools}{table.6}{main.pdf}}
